\section{Delay}
\begin{definition}[Delay model]
	The delay model an enhanced version of the Semi-Online model, where on vertex arrival it doesn't have to be matched immediately, but only after a delay of $f(n)$ time - $f(n)-1$ arriving online vertices, meaning that when $t+f(n)$-th vertex arrives, all vertices up to and including $t$-th must be irrevocably matched.
\end{definition}

In this way there is a delay window of $f(n)$ unmatched vertices, one of which has to be matched immediately and a vertex has to be irrevocably matched after $f(n)$ time has passed.

This model isn't well-studied as despite the fact that delay (sometimes referred to as buffer or lookahead) is often utilized, as it isn't useful in the case of Online Bipartite Matching

\begin{theorem}
	For delay no greater than $n/2$ no deterministic online algorithm has competitiveness greater than $0.5$
\end{theorem}

\begin{proof}
	Consider the following graph $G$: There are $k=n/2$ groups $\{Z_i|i\in[k]\}$. of two offline vertices\\
	First, $k$ vertices $z_i^D$ arrive, each connected to both vertices of $Z_i$.\\
	As $z_k^D$ arrives the cache is full, and $z_1^D$ has to be matched. Let the offline vertex it was matched to be $z_1^B$. Next, $z_1^C$ arrives, only connected to $z_1^B$, which is unmatchable. Let $z_1^A$ be the unmatched offline vertex.\\
	After $z_1^C$ has arrived $z_2^D$ has to be matched. The procedure above is repeated until $z_k^C$ has arrived.\\
	The perfect matching of $G$ connects all $z_i^A,z_i^D$ and $z_i^B,z_i^C$ pairs, while the algorithm connects only $z_i^B,z_i^D$ pairs, making the competitiveness $0.5$. The final graph is illustrated in the figure below.
\end{proof}

\begin{figure}[H]
	\centering
	\scalebox{.8}{\begin{tikzpicture}
	\begin{pgfonlayer}{nodelayer}
		\node [style=default] (0) at (-1, 4) {};
		\node [style=default] (1) at (1, 4) {};
		\node [style=default] (2) at (-1, 3) {};
		\node [style=default] (3) at (1, 3) {};
		\node [style=nil] (4) at (-1, 4.5) {$z_1^A$};
		\node [style=nil] (5) at (1, 4.5) {$z_1^D$};
		\node [style=nil] (6) at (-1, 3.5) {$z_1^B$};
		\node [style=nil] (7) at (1, 3.5) {$z_1^C$};
		\node [style=nil] (8) at (-1.5, 4) {};
		\node [style=nil] (9) at (-1.5, 3) {};
		\node [style=nil] (10) at (-2, 3.5) {$Z_1$};
		\node [style=default] (11) at (-1, 2) {};
		\node [style=default] (12) at (1, 2) {};
		\node [style=default] (13) at (-1, 1) {};
		\node [style=default] (14) at (1, 1) {};
		\node [style=nil] (15) at (-1.5, 2) {};
		\node [style=nil] (16) at (-1.5, 1) {};
		\node [style=default] (17) at (-1, -0.25) {};
		\node [style=default] (18) at (1, -0.25) {};
		\node [style=default] (19) at (-1, -1.25) {};
		\node [style=default] (20) at (1, -1.25) {};
		\node [style=nil] (21) at (-1.5, -0.25) {};
		\node [style=nil] (22) at (-1.5, -1.25) {};
		\node [style=nil] (23) at (-2, 1.5) {$Z_2$};
		\node [style=nil] (24) at (-2, -0.75) {$Z_k$};
		\node [style=nil] (25) at (0, 0.75) {.};
		\node [style=nil] (26) at (0, 0.5) {.};
		\node [style=nil] (27) at (0, 0.25) {.};
	\end{pgfonlayer}
	\begin{pgfonlayer}{edgelayer}
		\draw (0) to (1);
		\draw (2) to (3);
		\draw [bend right=90, looseness=0.75] (8.center) to (9.center);
		\draw (11) to (12);
		\draw (13) to (14);
		\draw [bend right=90, looseness=0.75] (15.center) to (16.center);
		\draw (17) to (18);
		\draw (19) to (20);
		\draw [bend right=90, looseness=0.75] (21.center) to (22.center);
		\draw [style=greenl] (1) to (2);
		\draw [style=greenl] (12) to (13);
		\draw [style=greenl] (18) to (19);
	\end{pgfonlayer}
\end{tikzpicture}
}
\end{figure}

\begin{theorem}
	For any $f(n)=o(n)$ delay, no algorithm has asymptotic competitiveness better than $1-\frac{1}{e}$
\end{theorem}
\begin{proof}
	Let $m=\ceil{f(n)}+1$.
	For $\mathrm{ALG}$ construct the adversary graph in a manner similar to the one above, with $m$ isolated $I^\star$ (from 3.2) adversarial subgraphs whose construction will be described after their purpose is clear, and vertices arriving to each subgraph in the same adversarial order while all subgraphs only have one online vertex in the delay window, implying that they receive their $k$-th online vertex before any receives it's $k+1$'st online vertex. Vertex labels are random.

	The behavior of the algorithm, isolated to a single subgraph, creates a family of some Semi-Online algorithms. This is evident as an algorithm for Semi-Online Matching can pretend that there are $m-1$ other subgraphs and this specific delay window exists. By solving this false instance, a true Semi-Online instance is also solved. To achive a competitiveness ratio above $1-\frac{1}{e}$, at least one of the subgraphs algorithms would have achived this same higher ratio on $I^\star$, which we know is impossible from section 3.2, where it is proven that, if labels are random, all algorithms have a performance upper bound of $1-\frac{1}{e}$ on $I^\star$.
	
	Because $f(n)=o(n)$ $\lim_{n\to\infty}\frac{n}{f(n)}=\infty$, implying that the subgraphs grow in size infinitely as the instance size grows, resulting in a competitiveness of $1-\frac{1}{e}$.
\end{proof}

\section{Corrections}
\autoref{thm:recourse} proves that there is an algorithm that maintains a maximum matching with amortized $O(\log^2n)$ corrections per arriving vertex. This produces a question: what would be the competitiveness of an algorithm that has a strict correction budget. The following model and proof are an extension of those ideas, producing an algorithm that doesn't guarantee a maximum matching but has a good competitiveness under more strict resource requirements.

\begin{definition}[Correction Model]
	In the correction model, every turn, the algorithm can remove up to a constant $d$ edges from the formed matching. There are no restrictions on adding edges
\end{definition}
An algorithm that improves the matching by changing some $d$ edges for each arrival, for some constant $d$, can be derived from a known lemma proven by Hopcroft and Karp in 1973 \cite{hopcroft-karp}.
\begin{lemma}
	for a matching $M$, if all augmenting paths are of length no less than $2k-1$, then $\frac{|M|}{|M^\star|}\geq\frac{k-1}{k}$
\end{lemma}
we will follow the proof from \cite{correction}
\begin{proof}
	Consider that graph $G'=(U,V,E')$ where $E'$ is the symmetric difference between some matching $M$ and $M^\star$, which is $M\triangle M'=(M\backslash M^\star)\cup(M^\star\backslash M)$, which in turn is the difference between $M$ and $M^\star$. Vertices in $G'$ have a degree of either $0,1$ or $2$. Vertices which are either isolated or in cycles (which must be of even length) in $G'$ have the same amount of edges adjacent from $M,M^\star$. Augmenting paths, of which there are no more than $|M^\star|-|M|$, must account for the difference in cardinality.

	The difference in cardinality between $M$ and $M^\star$ is the difference between the amount of their edges in the symmetric difference graph $G'$, which in turn is the sum over connected components, being paths and cycles. Therefore, the difference in cardinality is greatest, $G'$ contains only augmenting paths, of which each has at least $k-1$ edges in $M$ and $k$ edges in $M^\star$. Let $P$ be a path. Then $\frac{|M\cap P|}{|M^\star\cap P|}\geq\frac{k-1}{k}$ which can be generalized over all paths to $\frac{|M|}{|M^\star|}\geq\frac{k-1}{k}$
\end{proof}
From here we can construct the algorithm. We will call an augmenting path short if it is of length no greater than $2d+1$, meaning it has $d$ edges in the matching. Notice that an augmenting path of length $2k-1$ only has $k-1$ edges from the matching.
\begin{algorithm}
	\caption{$d$-augmenting correction}
	\label{alg:correction}
	\begin{algorithmic}
		\item on Arrival of vertex $v$: if short augmenting path's (starting at $v$) are created, switch the shortest. otherwise do nothing
	\end{algorithmic}
\end{algorithm}

\begin{theorem}
	$d$-augmenting correction has a competitiveness of $\frac{d+1}{d+2}$
\end{theorem}

We note that \cite{obmCorrections} proved the same competitiveness for this algorithm and it's optimality in a very similar model, but they are not identical so their optimality findings cannot be carried over without proof. Their competitiveness proof can be carried over directly to this model.

\begin{proof}
	We will prove a simple invariant on time $t$, that there are no short augmenting paths, ie shorter than $2d+3$ which directly implies the theorem. Though the previously mentioned model introduced by \cite{obmCorrections} is different, their approach for proving competitiveness can be used in this model. Do note that the 'shortest' prioritization in the algorithm is necesarry.\\
	At time $t=0$ there are no augmenting paths and the invariant holds.
	As $v$ arrives, assume that a shortest short $M$-augmenting path $P$ is created, the current matching is $M$, and $M$ augmented by $P$ is $M'$. Let $Q$ be any $M'$-augmenting path. We will prove that there exist $M$-augmenting paths $X,Y$ such that
	\begin{itemize}
		\item $X,Y$ are edge-disjoint
		\item $X\cup Y\subseteq P\cup Q$
		\item the sets of endpoints for $P,Q$ and $X,Y$ are identical
	\end{itemize}
	since all vertices in $P$ are matched in $M'$, trivially $Q$ has a separate set of endpoints from $P$. In a symmetric difference graph $H=(U,V,P\triangle Q)$ Four vertices, the endpoints of $P,Q$, have a degree of $1$ and all others have a degree of $2$ since otherwise that vertex would be internal to $Q$ and on $P$, implying the existance of an edge $e$ in $P\backslash M$ connected to it, implying that it is in $M'$, and in turn in $Q$, leading to a contradiction as this edge is also in $P\triangle Q$.\\
	With this structural fact we can conclude that $H$ contains two paths and some cycles. Let $X,Y$ be these paths. The properties of $X,Y$ are apparent, now we have to prove that they are augmenting paths.\\
	First, the endpoints of $Q$ must be unmatched in $M$ as they are not on $P$. notice that an edge in $Q\backslash P$ is in $M$ iff it is in $M'$. Let vertex $v$ have degre $2$ in $H$. If both of it's incident edges in $H$ are from $P$ or $Q$ then the statement holds. Else, $v$ in on $P$ and $Q$, and isn't and endpoint of $Q$. Similar to before, there must be an edge $e$ incident to $v$ that is $P\backslash M$ and $Q\cap M'$, implying that one of the incident edges from $H$, the one in $P\triangle Q$ is in $P\cap M$ and the other edge is in $Q\backslash(M\triangle P)=Q\triangle M$, making $X,Y$ augmenting paths.\\
	
	WLOG, suppose that the arriving $v$ is an endpoint of $X$. Trivially, it is longer than $P$ since they are both $M$-augmenting paths and the shorter on is used by the algorithm to augment $M$ to $M'$. Further, $v$ isn't an endpoint of $Y$, which is also an $M$-augmenting path. By induction hypothesis, both $X$ and $Y$ are not short. $X\cup Y\subseteq P\cup Q$ implies $|X|+|Y|\leq|P|+|Q|$, implying $|X|+|Y|\leq|X|+|Q|$ resulting in $|Y|\leq|Q|$, meaning that $Q$ must not be short.
\end{proof}